\documentclass[10pt]{article}

\usepackage[margin=1in]{geometry}
\usepackage{amsmath,amssymb,amsthm}
\usepackage{booktabs}
\usepackage{enumitem}
\usepackage{graphicx}
\usepackage[hidelinks]{hyperref}
\usepackage[numbers,sort&compress]{natbib}
\usepackage{xcolor}

\newtheorem{definition}{Definition}
\newtheorem{proposition}{Proposition}
\newtheorem{remark}{Remark}

\newcommand{\E}{\mathbb{E}}
\newcommand{\C}{\mathcal{C}}
\newcommand{\N}{\mathcal{N}}
\newcommand{\novel}{\star}
\newcommand{\method}{\textsc{Cira}}

\setlist{itemsep=1pt,topsep=3pt}

\title{World Models That Remember:\\ Adaptation as Contextual Inference}
\author{Research Proposal --- v1}
\date{September 28, 2026}

\begin{document}
\maketitle

\begin{abstract}
Test-time adaptation of latent world models treats every distribution shift as new: a few gradient
steps on a self-supervised prediction loss, then a reset at the next episode. A deployed robot does
not face new shifts; it faces the same few surfaces, payloads and actuator conditions again and
again, usually with something visible that correlates with them. This proposal treats deployment-time
adaptation as \emph{contextual inference}: a single posterior over a growing set of dynamics
contexts decides when a memory is created, which memory is updated, and how strongly each is
expressed. It follows the COIN account of human motor learning~\cite{heald2021coin}, in which part of
what looks like learning is not parameter change but a change in which existing memory is being
expressed. We turn that distinction into a measurement: recovery after a shift splits exactly into
\emph{apparent} learning (recall and re-expression) and \emph{proper} learning (parameter updates), and
per-episode-reset methods forgo the first term by construction. The method, \method{}
(Context-Indexed Residual Adapters), attaches closed-form Bayesian residual adapters to any frozen
latent world model, runs a nonparametric multiple-model filter over them, and learns during deployment,
without labels, which pre-contact visual features predict which context. The central, falsifiable
prediction is that, on recurring conditions with an informative cue, prediction error is corrected
\emph{before the first contact}, and that adaptation latency falls with each re-exposure by an amount a
simple sequential-testing argument predicts in advance. Evaluation uses long regime-switching streams
without oracle resets, built around four motor-adaptation paradigms, on a toy pusher, the AdaJEPA
environments~\cite{wang2026adajepa}, and real-robot pushing.
\end{abstract}

%=====================================================================
\section{Thesis and contributions}
%=====================================================================

Latent world models trained on large offline datasets now support zero-shot planning in manipulation
and navigation~\cite{zhou2025dinowm,assran2025vjepa2,sobal2025pldm,hansen2024tdmpc2}. Their
failure mode under deployment shift has been addressed very recently by self-supervised test-time
adaptation: AdaJEPA takes gradient steps on recent latent prediction errors inside the MPC
loop~\cite{wang2026adajepa}, and Sandwich-Residuals learns small residual modules around a frozen
predictor~\cite{soni2026sandwich}. Both reset to the pretrained model at every episode. Both name
persistence across episodes, or combination with continual and active learning, as future work.

Resetting is not a design oversight; it is what keeps gradient-based test-time adaptation stable. On
long streams, persistent updates drift and collapse, and periodic resets beat most continual TTA
methods~\cite{press2023rdumb,hoang2024petta}. The price is that every recurrence of a known condition
is paid for again, from scratch, with prediction error.

\paragraph{Thesis.}
Deployment-time adaptation of a world model should be treated as inference over \emph{which
dynamics context is active}, not as optimisation of one set of adapter parameters. Replacing ``update
the adapter'' with ``infer the context, then update only that context's adapter'' has three
consequences that gradient-based adaptation cannot reproduce: recurring conditions are recalled rather
than relearned (\emph{savings}), learning one condition does not overwrite another (\emph{no
interference}), and a visual cue that has been associated with a context shifts the posterior before
any prediction error arrives (\emph{anticipation}).

\paragraph{Contributions.}
\begin{enumerate}[label=\textbf{C\arabic*.},leftmargin=2.2em]
  \item \textbf{A measurement.} An exact two-player Shapley decomposition of post-shift recovery into
  apparent learning (changes in context expression) and proper learning (changes in adapter
  parameters), computable for any method by counterfactual replay (Def.~\ref{def:decomp}). We report,
  on realistic streams, how much of the recovery each term accounts for. The apparent term is
  identically zero for every per-episode-reset method.
  \item \textbf{A method.} \method{}: context-indexed Bayesian residual adapters in a low-dimensional
  error subspace, one nonparametric context posterior that drives creation, recall, reversion and
  change detection, a state-dependent noise model so that contact transients are not read as shifts,
  and cue likelihoods learned during deployment from the context's own dynamics evidence. It plugs into
  any frozen latent world model and has closed-form updates.
  \item \textbf{A quantitative prediction.} A sequential-testing argument (Prop.~\ref{prop:latency})
  predicts adaptation latency as a function of cue reliability, exposure count and the dynamics
  divergence between contexts. The neuroscience paradigms are therefore tested against numbers
  written down before the experiment, not against qualitative resemblance.
  \item \textbf{A protocol.} Long, stochastic, regime-switching streams without oracle resets, with the
  paradigms savings, anterograde interference, evoked recovery and cue conflict embedded as
  sub-streams, and a cue-reliability sweep. Released with the toy and AdaJEPA-environment variants.
  \item \textbf{A real-robot study.} Planar pushing on swapped surfaces whose texture correlates with
  friction, and payloads in visually distinct containers, with cue reliability controlled by the
  experimenter.
\end{enumerate}

\paragraph{Scope.}
Dynamics shifts (friction, mass, payload, actuator gain) under a frozen encoder. Appearance enters only
as evidence about the context, never as something to correct; appearance shifts that degrade the
encoder itself are the subject of AdaJEPA and Sandwich-Residuals and are out of scope.

%=====================================================================
\section{Background and positioning}
%=====================================================================

\paragraph{Contextual inference in motor learning.}
MOSAIC proposed a repertoire of paired forward and inverse models whose responsibilities are set by
both prediction accuracy and contextual signals available before the movement
starts~\cite{wolpert1998multiple,haruno2001mosaic}. COIN~\cite{heald2021coin} made this a
generative model: an unbounded set of contexts with sticky Markov transitions under a hierarchical
Dirichlet process, a linear-Gaussian memory per context, and context-dependent sensory cues, inverted
by particle learning. COIN explains savings, anterograde interference and spontaneous recovery through
changes in predicted context probabilities rather than changes in learning rate. It also predicted and
confirmed \emph{evoked recovery}, which the two-state model~\cite{smith2006interacting} cannot
produce, and showed that single-trial learning is graded by the combined evidence of cue and feedback.
Reviews extend the framework across memory
systems~\cite{heald2023tics,heald2023annurev,shivkumar2025curriculum}. The same logic appears in
latent-cause models of conditioning, where gradual change updates an existing memory and abrupt change
creates a new one~\cite{gershman2010context,gershman2014statistical}, and in error-driven memory
switching~\cite{oh2019minimizing}. Human data also give a caution we take seriously: many static visual
cues are weak at separating dynamics memories~\cite{howard2013contextual}, while cues tied to the
movement plan are strong~\cite{sheahan2016planning}. A cue has to earn its association, which is why
\method{} learns cue likelihoods from dynamics evidence and why we sweep cue reliability rather than
assume it.

\paragraph{Test-time adaptation of world models.}
AdaJEPA~\cite{wang2026adajepa} and Sandwich-Residuals~\cite{soni2026sandwich} are the direct
predecessors. AdaJEPA adapts the last encoder and selected predictor layers with about one gradient step
per replan on the five most recent transitions, evaluated on PushT, PushObj and PointMaze under
visual, shape, dynamics (mass, damping) and layout shifts. It states that adaptation is bounded by the
coverage of the pretrained representation. Sandwich-Residuals reaches comparable gains with around
1--3\% of the parameters by adapting only residual modules. Neither carries anything across episodes,
represents uncertainty, or uses exteroceptive evidence proactively. General TTA
methods~\cite{sun2020ttt,wang2021tent,wang2022cotta} share the same template, and their failure on
long and recurring streams~\cite{press2023rdumb,hoang2024petta} is the empirical motivation for
memory. Recent world-model work detects shifts and forgets stale data~\cite{yang2026cawm} or infers a
continuous context reactively~\cite{roder2025dali,benad2026dmash}, but does not store and recall
discrete contexts.

\paragraph{Bayesian context machinery for dynamics.}
The components of \method{} are individually established. Bayesian last layers with closed-form online
updates~\cite{harrison2018alpaca,harrison2024vbll}; changepoint-driven switching of such
models~\cite{harrison2020moca,adams2007bocpd}; CRP or DP mixtures of dynamics models learned
online~\cite{nagabandi2019mole,xu2020taskagnostic,jerfel2019reconciling,lee2020cndpm}; switching
linear dynamical systems with nonparametric priors~\cite{fox2011bnpslds,fox2011sticky,linderman2017rslds};
and, in estimation, the interacting multiple-model filter~\cite{blom1988imm,mazor1998imm} with
variable model sets~\cite{li1996vsmm}. MOLe~\cite{nagabandi2019mole} is the closest ML precedent: an
EM-fitted CRP mixture of meta-learned dynamics models for model-based RL. It selects contexts purely
by prediction error, so it cannot anticipate. MOCA~\cite{harrison2020moca} discards the old context at
every changepoint, so it cannot recall. None of these operate on the latent space of a pretrained
visual world model.

\paragraph{Vision-informed proactive dynamics.}
Vision-Conditioned VBLL~\cite{brunzema2026vcvbll} argues that reactive adaptation is too slow and
conditions a Bayesian last-layer vehicle model on a hand-designed water score, driving through
puddles that non-visual adaptive baselines could not handle. Terrain-aware
dynamics~\cite{gibson2024visualterrain,cai2023probtrav} and vision-based
locomotion~\cite{agarwal2022egocentric} make the same point. Closest in spirit to deployment-learned
cues, cross-modal supervision trains vision online to predict proprioceptively inferred
quantities~\cite{loquercio2022crossmodal,margolis2023activesensing,frey2023wvn}. These methods map
appearance to a continuous physical quantity through a pretrained or pre-specified channel. \method{}
instead maps appearance to a \emph{discrete, self-discovered} context, with no labels and no training
on the shift, and lets the dynamics evidence override the cue.

\begin{table}[t]
\centering
\footnotesize
\setlength{\tabcolsep}{2.5pt}
\resizebox{\linewidth}{!}{%
\begin{tabular}{lcccccc}
\toprule
 & Latent visual & Persists across & Recalls & Calibrated & Pre-contact & Cue learned \\
 & world model & episodes & contexts & uncertainty & anticipation & in deployment \\
\midrule
AdaJEPA~\cite{wang2026adajepa}, Sandwich-Res.~\cite{soni2026sandwich} & \checkmark & -- & -- & -- & -- & -- \\
MOLe~\cite{nagabandi2019mole}, IMGP~\cite{xu2020taskagnostic} & -- & \checkmark & \checkmark & partial & -- & -- \\
MOCA~\cite{harrison2020moca} & -- & \checkmark & -- & \checkmark & -- & -- \\
VcVBLL~\cite{brunzema2026vcvbll} & -- & -- & -- & \checkmark & \checkmark & -- (hand-designed) \\
Cross-modal vision~\cite{loquercio2022crossmodal} & -- & \checkmark & -- & -- & \checkmark & \checkmark (continuous) \\
COIN~\cite{heald2021coin} (human model) & -- & \checkmark & \checkmark & \checkmark & \checkmark & \checkmark \\
\midrule
\method{} (this proposal) & \checkmark & \checkmark & \checkmark & \checkmark & \checkmark & \checkmark \\
\bottomrule
\end{tabular}}
\caption{Positioning. As of September 2026 we found no work combining a frozen latent visual world
model, a nonparametric posterior over stored dynamics contexts, and cues learned during deployment.
Each column is individually covered; the combination and the measurement (C1) are the contribution.}
\label{tab:positioning}
\end{table}

%=====================================================================
\section{Method: Context-Indexed Residual Adapters}
%=====================================================================

\paragraph{Setup.}
A frozen encoder $E$ maps observations to latents $z_t=E(o_t)\in\mathbb{R}^d$, and a frozen predictor
gives $\hat z_{t+1}=f(z_t,a_t)$. The world model is used by a sampling-based planner (CEM or
MPPI)~\cite{rubinstein1999cem,chua2018pets}. Deployment is a single stream, with known episode
boundaries but no regime labels and no reset of any adaptive state.

\paragraph{Residual subspace.}
On nominal (pretraining or pre-deployment) data, collect residuals $r_t=z_{t+1}-f(z_t,a_t)$ and let
$U\in\mathbb{R}^{d\times k}$ hold their top-$k$ principal directions, $k\approx 4$--$16$. All adaptation
happens in $e_t=U^\top r_t\in\mathbb{R}^k$. This is the hidden-parameter view of dynamics
shifts~\cite{doshivelez2016hipmdp}: a few physical parameters move the dynamics along a few
directions. The subspace is a bet that these directions coincide with where the nominal model already
errs. Section~\ref{sec:risks} gives an online diagnostic for when the bet fails.

\paragraph{Context-indexed Bayesian adapters (the memories).}
Each context $c$ holds a Bayesian linear residual model on fixed features
$\phi_t=\phi(U^\top z_t, a_t)\in\mathbb{R}^p$, for example $[a_t,\;U^\top z_t,\;1]$:
\begin{equation}
  e_t \mid c \;\sim\; \N\!\big(W_c\,\phi_t,\; \Sigma(z_t,a_t)\big),\qquad
  W_c \sim \mathcal{MN}(M_c, \Sigma_{\text{row}}, V_c).
  \label{eq:adapter}
\end{equation}
This departs from the working draft, which used a Gaussian bias. A friction or gain change produces a
residual that scales with the action, so a bias cannot represent it; a linear-in-features residual can,
and it keeps closed-form conjugate updates as in ALPaCA~\cite{harrison2018alpaca}. Within a context,
$W_c$ follows a slow random walk (process noise $Q$). This gives each memory its own proper learning
and mild forgetting, the analogue of COIN's per-context retention and drift. The \emph{base} context
$c=0$ has $W_0\equiv 0$ and reproduces the frozen model. The \emph{novel} context $c=\novel$ carries
the broad prior $\mathcal{MN}(0,\Sigma_{\text{row}},V_0)$.

\paragraph{State-dependent noise.}
$\Sigma(z,a)$ is a small heteroscedastic head fitted on nominal residuals by Gaussian negative
log-likelihood~\cite{kersting2007hetgp}, shared across contexts. Contact onsets, stick--slip and
impacts inflate $\Sigma$ exactly where nominal prediction is already poor, which shrinks the
log-likelihood ratio between contexts there. Planar pushing is known to be heteroscedastic and
friction-variable~\cite{bauza2017probpush,ma2018friction}. We found no prior work that makes the link
between contact transients and false context creation explicit~\cite{haddadin2017collisions}; we
expect it to be the most likely failure of a naive implementation, so the toy includes it deliberately.

\paragraph{One posterior over contexts.}
Let $\C_t=\{0,1,\dots,K_t\}\cup\{\novel\}$. The prior over the context of a new episode is a
CRP-style predictive with a sticky continuation, and within an episode a sticky Markov
chain~\cite{fox2011sticky}:
\begin{equation}
  p(c_t=j\mid c_{t-1}=i) \propto
  \begin{cases}
    \kappa+\alpha_{\text{sw}}\, n_j & j=i\\
    \alpha_{\text{sw}}\, n_j & j\neq i,\; j\in\{0,\dots,K_t\}\\
    \alpha_{\text{sw}}\,\gamma & j=\novel,
  \end{cases}
  \label{eq:prior}
\end{equation}
where $n_j$ is the soft occupancy count of context $j$, $\gamma$ the concentration, and $\kappa$ the
self-transition bias. The filter is an IMM~\cite{blom1988imm} over the adapters in
Eq.~\eqref{eq:adapter}. Each context's predictive for $e_t$ is Student-$t$ (or Gaussian after
marginalising $W_c$), and the posterior $\pi_t(c)\propto p(e_t\mid c,\phi_t)\,p(y_t\mid c)\,
\sum_i p(c\mid i)\,\pi_{t-1}(i)$ combines dynamics and cue evidence. Each adapter is then updated with
weight $\pi_t(c)$, a responsibility-weighted conjugate update, not winner-take-all, as in
COIN~\cite{heald2021coin}. The one posterior serves four roles, with no separate detector:
\begin{itemize}
  \item \emph{creation}: the novel slot dominates over a short window, so its posterior is cloned into
  a new stored context and the slot is reset to the prior;
  \item \emph{recall}: a stored context dominates;
  \item \emph{reversion}: the base model dominates again;
  \item \emph{change detection}: the entropy or MAP of $\pi_t$ changes.
\end{itemize}
Contexts whose parameter posteriors come within a small symmetric KL of each other are merged, which
bounds $K_t$ and repairs duplicated contexts.

\paragraph{Deployment-learned cues.}
At the start of each episode, before first contact, the robot observes $y=P^\top E(o_{\text{pre}})$, a
low-dimensional projection ($\le 8$ dims, principal components of nominal pre-contact latents) of the
scene. Each context holds a conjugate Normal--Wishart cue likelihood $p(y\mid c)$. The design choice
that matters is that \emph{cue likelihoods are updated with responsibilities computed from dynamics
evidence only}, smoothed over the finished episode, never from the cue-augmented posterior. This
prevents the self-confirming loop in which a cue assigns a context and that assignment then reinforces
the cue. It also means a misleading cue can be unlearned, which H4 tests. Once associated, the cue term
$p(y\mid c)$ shifts $\pi$ at $t=0$, and the planner uses the posterior-weighted adapter before any
contact.

\paragraph{Planning.}
The planner rolls out $\hat z_{t+1}=f(z_t,a_t)+U\,\bar W\phi_t$ with $\bar W=\sum_c\pi_t(c)M_c$ (mixture
mean), or samples contexts per CEM particle for risk-sensitive costs. All per-step filtering is
$O(|\C_t|\,k\,p^2)$ with $k,p\lesssim 16$, negligible next to a single predictor call.

\paragraph{Active disambiguation (stretch).}
When $H(\pi_t)$ exceeds a threshold, CEM ranks elite candidates by
$J(\mathbf a)=J_{\text{task}}(\mathbf a)-\beta\, I(c;\,e_{t+1:t+H}\mid \mathbf a)$. The mutual information
is between the discrete context and a Gaussian mixture of predicted residuals, and we bound it in closed
form with pairwise divergences between the per-context predictives. The task-cost bound is enforced by
construction: only candidates with $J_{\text{task}}\le(1+\epsilon)\min J_{\text{task}}$ are eligible for
the information term. This is active fault diagnosis~\cite{blackmore2006discrimination,
blackmore2008active,simandl2009afd} and dual control~\cite{mesbah2018dual,heirung2019mpcdiscrimination}
moved into latent planning. It differs from OPAX~\cite{sukhija2023opax} in targeting information about
a discrete context rather than about model parameters.

%=====================================================================
\section{The measurement: apparent versus proper learning}
%=====================================================================

Every adaptive world model in this proposal, including the baselines, has an adaptive state that
splits into an \emph{expression} part $\pi$ and a \emph{parameter} part $\Theta$. For \method{} these
are the context posterior and the adapter posteriors. For MOLe they are the mixture weights and the
network weights. For AdaJEPA and SGD-based adapters, $\pi$ is trivial and everything is $\Theta$. Let
$L(\pi,\Theta)$ be the one-step latent prediction error of the resulting model on transitions from the
current regime, evaluated by replay.

\begin{definition}[Recovery decomposition]\label{def:decomp}
For a shift at time $t_0$, the recovery by time $t$ is
$R_t=L(\pi_{t_0},\Theta_{t_0})-L(\pi_t,\Theta_t)$. Its apparent and proper parts are the two-player
Shapley values
\begin{align*}
  A_t &= \tfrac12\big[L(\pi_{t_0},\Theta_{t_0})-L(\pi_t,\Theta_{t_0})\big]
        +\tfrac12\big[L(\pi_{t_0},\Theta_t)-L(\pi_t,\Theta_t)\big],\\
  P_t &= R_t-A_t,
\end{align*}
with contexts created after $t_0$ given zero mass under $\pi_{t_0}$.
\end{definition}

The decomposition is exact ($A_t+P_t=R_t$) and order-invariant. It is cheap for \method{} because every
counterfactual is a closed-form evaluation of stored posteriors. It makes the COIN distinction
operational: on a first exposure recovery should be almost entirely proper; on a re-exposure it should
be almost entirely apparent. For any method that resets at episode boundaries, or has no expression
state, $A_t\equiv 0$. We report $A_t/R_t$ over realistic streams as a headline number, not only on the
paradigm sub-streams.

%=====================================================================
\section{A quantitative prediction for latency, savings and cues}
%=====================================================================

The paradigms earn their place only if they separate methods quantitatively. The following
argument turns them into numbers fixed before the experiment.

\begin{proposition}[Latency as a first-passage time]\label{prop:latency}
Consider an episode whose true context $c^\ast$ is stored, and treat the strongest alternative $c'$ as
the competing hypothesis. Let the per-step dynamics log-likelihood ratios $\ell_s$ after first contact
$t_c$ be i.i.d.\ with mean $D=\E_{c^\ast}[\mathrm{KL}]>0$, the expected divergence between the two
contexts' residual predictives under the state-dependent noise $\Sigma$. Before contact,
$\ell_s=0$. Let the episode-start log-odds be
\[
  \lambda_0 \;=\; \underbrace{\log\frac{n_{c^\ast}}{n_{c'}}}_{\text{prior (CRP counts)}}
  \;+\;\underbrace{\log\frac{p(y\mid c^\ast)}{p(y\mid c')}}_{\text{cue, }\ell_y},
\]
and let latency $\tau$ be the number of post-contact steps until the log-odds first reach
$\lambda^\ast$. Then $\tau=0$ if $\lambda_0\ge\lambda^\ast$, and otherwise, by Wald's identity,
$\E[\tau]=(\lambda^\ast-\lambda_0+\E[\text{overshoot}])/D\ge(\lambda^\ast-\lambda_0)/D$.
\end{proposition}

The proposition is elementary. Its value is in what it commits us to.
\begin{itemize}
  \item \textbf{Anticipation (H1).} Latency is zero exactly when prior plus cue clear the threshold.
  With soft expression, the fraction of the shift corrected before contact is $\pi_0(c^\ast)$, which
  for a categorical cue of reliability $\rho$ over $K$ equiprobable contexts is $\rho$ once the cue is
  learned. With the toy's $\rho=0.8$, about 80\% of the shift should be corrected at $t_c$.
  \item \textbf{Savings (H2).} $\lambda_0$ grows like $\log n_{c^\ast}$, so latency falls with each
  re-exposure by $\approx\log(n_{k+1}/n_k)/D$ steps, with diminishing returns. For reset-based SGD,
  latency is constant in $k$.
  \item \textbf{Cue conflict (H4).} Relative to a correct cue, a misleading cue of strength $\ell_y$ costs
  $\approx 2\ell_y/D$ extra steps before the dynamics override it. The override is bounded, and it
  scales inversely with how different the contexts' dynamics are.
  \item \textbf{Noise (H5).} Heteroscedastic $\Sigma$ lowers $D$ at contact transients. That costs
  latency but prevents the false contexts a homoscedastic model would create. We measure both sides of
  this trade-off.
\end{itemize}
We will fit $D$ from nominal data, predict latency curves, and report prediction against observation.
A large, systematic mismatch is itself a finding, most likely non-i.i.d.\ evidence around contact
events.

%=====================================================================
\section{Hypotheses and kill criteria}
%=====================================================================

\begin{enumerate}[label=\textbf{H\arabic*},leftmargin=2.2em]
  \item \textbf{Anticipation.} On recurring conditions with an informative cue, prediction error at
  first contact is reduced by roughly the posterior mass on the true context, and latency is $\le 0$.
  The cue-free bank needs $\approx(\lambda^\ast-\lambda_0)/D>0$ steps.
  \item \textbf{Savings.} Recovery time drops with each re-exposure, following
  Prop.~\ref{prop:latency}. For SGD and SGD with periodic reset it stays roughly constant.
  \item \textbf{No interference.} Learning regime B leaves prediction and task performance on A
  intact. Always-on SGD degrades A, as do continual-TTA baselines.
  \item \textbf{Cue conflict.} With a misleading cue, dynamics evidence overrides the cue within the
  bounded number of steps predicted above, and task cost stays below the frozen model's.
  \item \textbf{Noise robustness.} Under stochastic, contact-rich dynamics, no single SGD learning rate
  is best in both the noisy-nominal and post-shift phases, while \method{} is no worse than the best
  one in either. \method{}'s false-context rate on nominal streams stays below a pre-registered bound.
  \item \textbf{Decomposition.} On realistic streams with recurrence, a substantial share of recovery
  (pre-registered: $A/R\ge 0.3$ at the median shift) is apparent learning.
\end{enumerate}

\paragraph{Kill criteria.}
\begin{itemize}
  \item If the cue-free filtered bank matches the full model on recurrence speed, we drop the
  anticipation claim, and the paper becomes ``filtered bank plus protocol and measurement''.
  \item If tuned SGD with periodic reset matches the filtered bank, the paper becomes the protocol, the
  decomposition, and the negative finding.
  \item If $A/R$ is small on realistic streams, the reframing is weaker than claimed and we say so. The
  measurement remains a contribution.
\end{itemize}

%=====================================================================
\section{Evaluation}
%=====================================================================

\paragraph{Streams.}
A regime process over $\{$nominal, $R_1,\dots,R_M\}$ with geometric dwell times (in episodes) and
occasional within-episode switches. Each regime is paired with an appearance, and the cue shows the
paired appearance with probability $\rho$ and a random other one otherwise. We sweep
$\rho\in\{0.5,0.65,0.8,0.95\}$, with $\rho=1/M$ as the uninformative control. Dynamics are stochastic
throughout. Four paradigm sub-streams are embedded:
\begin{itemize}
  \item \textbf{Savings}: A--B--A--B blocks.
  \item \textbf{Anterograde interference}: a long versus short A block before B.
  \item \textbf{Evoked recovery}: after A then B, two A episodes are inserted into a nominal block; we
  test whether the A memory re-expresses.
  \item \textbf{Cue conflict}: cue pairings are swapped for one block.
\end{itemize}

\paragraph{Baselines.}
\begin{itemize}
  \item Frozen world model.
  \item Always-on SGD adapter (AdaJEPA-style).
  \item SGD with periodic reset (the RDumb lesson~\cite{press2023rdumb}).
  \item SGD with oracle reset at regime changes.
  \item Sandwich-Residuals~\cite{soni2026sandwich}.
  \item A MOLe-style CRP mixture of SGD adapters~\cite{nagabandi2019mole}.
  \item MOCA-style changepoint switching without recall~\cite{harrison2020moca}.
  \item \method{} ablations: cue-free, homoscedastic noise, bias-only adapter, no merging, and
  winner-take-all updates.
  \item Oracle-context \method{} as an upper bound.
\end{itemize}
All baselines see the same subspace and planner, and learning rates and filter hyperparameters are tuned
on nominal streams only.

\paragraph{Metrics.}
\begin{itemize}
  \item Adaptation latency relative to first contact.
  \item Excess prediction error integrated after each shift.
  \item Task cost and success.
  \item A savings ratio across exposures.
  \item Retention of A after B.
  \item The decomposition $A/R$.
  \item False-context and missed-context rates.
  \item Predicted-versus-observed latency (Prop.~\ref{prop:latency}).
\end{itemize}
Seeds and streams are pre-registered, and we report bootstrap confidence intervals.

\paragraph{Environments.}
\begin{enumerate}[label=(\roman*)]
  \item \textbf{Toy pusher.} A 2D pusher with contact-dependent noise, friction and mass regimes, and an
  80\%-reliable surface-colour cue, run on one savings and one cue-conflict stream. This is the go/no-go
  test for H1, H2, H4 and the noise risk, and it runs in minutes on a CPU.
  \item \textbf{AdaJEPA environments}~\cite{wang2026adajepa}: PushT, PushObj and PointMaze with their
  released JEPA and DINO-WM~\cite{zhou2025dinowm} checkpoints. PointMaze already has mass and damping
  shifts. We repurpose the existing appearance changes (recolouring the agent or block) as cues by
  pairing each with a dynamics shift, and add stochastic dynamics. This gives direct comparability with
  AdaJEPA and Sandwich-Residuals.
  \item \textbf{Real arm.} Planar pushing on swapped mats (e.g.\ felt, acrylic, rubber) whose texture
  correlates with friction~\cite{yu2016million}, and payloads in visually distinct containers. Cue
  reliability is controlled by occasionally covering mats with a neutral sheet or swapping container
  contents. The world model is a DINO-WM-style model trained on nominal pushing data.
\end{enumerate}

%=====================================================================
\section{Risks and mitigations}\label{sec:risks}
%=====================================================================

\begin{itemize}
  \item \textbf{Contact transients read as shifts.} Most likely failure. Mitigated by the
  heteroscedastic $\Sigma$ and a minimum dwell before creation, with the false-context rate reported
  as a first-class metric.
  \item \textbf{Weak cue--dynamics correlation.} Human data show many static visual cues are
  weak~\cite{howard2013contextual}. We choose tasks where the correlation is natural and sweep $\rho$
  rather than assume it. Prop.~\ref{prop:latency} says how much a weak cue should still help.
  \item \textbf{Spurious contexts from irrelevant appearance.} The cue enters only as a likelihood on a
  low-dimensional projection, trained on dynamics-only responsibilities, so appearance cannot
  \emph{create} a context. The CRP concentration is tuned on nominal streams.
  \item \textbf{Shift outside the error subspace.} Diagnosed online by the fraction of residual energy
  outside $U$. If it rises persistently, the novel context gets an extended subspace (the top
  directions of the unexplained residual). This is also where \method{} meets the coverage limit
  AdaJEPA reports.
  \item \textbf{Neuroscience window dressing.} Each paradigm is tied to a quantitative gap between
  methods and to a prediction from Prop.~\ref{prop:latency}. A paradigm that does not separate methods
  in the toy is dropped from the paper.
  \item \textbf{Scooping.} AdaJEPA and Sandwich-Residuals both name persistence as next steps, so a
  workshop version by spring matters. VcVBLL~\cite{brunzema2026vcvbll} already argues for proactive
  exteroceptive adaptation. Our differentiators are memory creation and recall, deployment-learned
  cues, the decomposition, and manipulation with latent world models.
\end{itemize}

%=====================================================================
\section{Plan}
%=====================================================================

\begin{center}
\small
\begin{tabular}{lp{0.42\linewidth}p{0.3\linewidth}}
\toprule
When & Work & Exit condition \\
\midrule
Oct 2026 & Toy pusher, \method{} core, SGD/reset baselines & H1/H2/H4 visible; kill criteria checked \\
Nov 2026 & Decomposition tooling; latency predictions vs.\ toy & Prop.~\ref{prop:latency} within error bars or understood \\
Dec 2026--Jan 2027 & AdaJEPA environments, full baseline set & Paradigm streams in PushT/PointMaze \\
Feb 2027 & Workshop paper (toy + AdaJEPA envs) & Submitted to a spring workshop \\
Mar--Apr 2027 & Real-arm pushing and payload study; active disambiguation & H1--H5 on hardware \\
May 2027 & Full paper & NeurIPS or CoRL 2027 submission \\
\bottomrule
\end{tabular}
\end{center}

\paragraph{What changed from the working draft.}
\begin{enumerate}
  \item The Gaussian bias adapter became a Bayesian linear residual, because friction and gain shifts are
  action-dependent.
  \item Cue likelihoods are learned from dynamics-only responsibilities, to prevent self-confirming cue
  loops.
  \item Apparent-versus-proper learning became an exact Shapley decomposition (C1).
  \item The paradigms are now tied to a sequential-testing latency prediction (C3), which answers the
  window-dressing risk.
  \item Evoked recovery and a decomposition hypothesis (H6) were added.
  \item Sandwich-Residuals (September 2026), MOLe and MOCA were added as baselines, since they are the
  nearest competitors found in the literature search.
  \item The active-disambiguation task-cost bound is enforced by construction through elite-set
  eligibility.
\end{enumerate}

\begingroup
\small
\setlength{\bibsep}{1pt}
\bibliographystyle{plainnat}
\bibliography{references}
\endgroup

\end{document}
